\section{Representative case studies}
\label{sec:case-studies}

\subsection{Selection policy and source provenance}

Examples can make an audit readable, but selected examples can also conceal a
failure mode or overstate a method's scope. The three cases in this section
are therefore not hand-picked for their visual contrast or effect magnitude.
The display configuration selects one complete supported-candidate anchor and
one complete not-supported anchor by minimum distance to the within-tier median
standardized score, with anchor-ID tie-breaking. It selects the
lexicographically first donor-insufficient inconclusive anchor to make the
absence of a qualified counterfactual visible. The configuration, selected
donors, archive manifest checksum, donor-table checksum, fitted weights, and
reconstruction error are stored in the generated case-study manifest.

For the two complete cases, the renderer reads only the checksum-pinned public
EPA archives, applies the frozen cleaning rule, reconstructs the fixed-prior
MetaShift weights from pre-anchor data, and fails if the reconstructed saved
log effect differs by more than $10^{-9}$. The raw archives remain outside
version control. The third case intentionally has no reconstructed
counterfactual: its audit disposition is that fewer than three qualified
distinct physical donors were available.

\begin{table}[tbp]
\centering
\caption{Deterministically selected representative audit cases.}
\label{tab:case-studies}
\footnotesize
\begin{tabular}{@{}p{0.16\linewidth}p{0.23\linewidth}p{0.17\linewidth}p{0.34\linewidth}@{}}
\toprule
Case & Metadata anchor & Disposition & Saved diagnostic summary \\
\midrule
Supported candidate & \texttt{12-057-0113}\newline\texttt{poc1-2023-12-01} & Complete & Effect -0.12777; fixed [-0.17258, -0.10511]; nested [-0.17399, -0.10560] \\
Not supported & \texttt{06-031-0004}\newline\texttt{poc8-2021-01-12} & Complete & Effect -0.08240; fixed [-0.21609, 0.03012]; nested [-0.24309, 0.03662] \\
Inconclusive & \texttt{01-103-0011}\newline\texttt{poc3-2023-02-01} & Donor insufficient & No counterfactual, effect, or interval is imputed. \\
\bottomrule
\end{tabular}
\par\smallskip
\begin{minipage}{0.94\linewidth}
\footnotesize\textit{Note.} Supported and not-supported cases minimize absolute distance to their within-tier median standardized score, with anchor ID tie-breaking. The inconclusive case is the lexicographically first donor-insufficient anchor and deliberately has no counterfactual or interval. Intervals are diagnostic, not calibrated physical-bias confidence intervals.
\end{minipage}
\end{table}


\begin{figure}[tbp]
\centering
\includegraphics[width=\linewidth]{fig_case_studies.pdf}
\caption{Deterministically selected representative cases. Top row: retained
daily target values and, for complete cases only, the pre-anchor-fitted
fixed-prior MetaShift counterfactual in $\mu$g/m$^3$. Middle row: calibrated
log residuals for complete cases; the inconclusive case instead states why no
counterfactual is constructed. Bottom row: protocol-defined fixed and nested
diagnostic intervals for complete cases; no effect or interval is imputed for
the donor-insufficient anchor. Blue and yellow shading mark the 60-day pre-
and post-anchor comparison windows. All raw display inputs are checksum-pinned
locally; the figure is not evidence of a physical instrument event.}
\label{fig:case-studies}
\end{figure}

\subsection{Interpretation boundaries}

The supported-candidate panel illustrates a case in which the predeclared
diagnostic conditions jointly hold. That status prioritizes the anchor for
further station-history and technical-record review; it does not identify a
device, a maintenance action, a calibration change, or a measurement bias. The
not-supported panel illustrates the opposite distinction: a visible target
series and a complete counterfactual are not sufficient when available
diagnostics fail the frozen evidence rules. It would be misleading to discard
that case merely because its residual plot is less dramatic.

The inconclusive panel is equally important. A target series may have an
apparent before/after pattern, yet the protocol cannot form the common
cross-site comparison if the distinct-donor requirement fails. MetaShift-Bench
therefore abstains rather than borrowing a same-site POC, relaxing the donor
rule after inspection, or assigning a physical explanation. This case makes
the audit's coverage boundary concrete: lack of a qualified comparison is
lack of evidence for this protocol, not evidence that nothing happened.

\FloatBarrier
